% SOURCE_AUTHORITY: sga5_fr_workpass.tex, lines 6483--6547 (LF-normalized unit).
% SOURCE_SHA256: 791F4EFFC5E02832D5D77ED03518C8156D6F07E4C8238B03545DB93D883FBB28
\begin{statement}{Proposición 6.23}
Bajo las hipótesis de 6.21, supongamos dada además una inmersión abierta
\[
i:V\hookrightarrow Y
\]
tal que $(fc_1)^{-1}(V)$ sea conexo, que $d_1^{-1}(V)\subset d_2^{-1}(V)$, y que $f$ induzca un recubrimiento principal de Galois de grupo $G$,
\[
f_V:U=f^{-1}(V)\to V.
\]
Supongamos además que $M$ es de la forma $M=i_!E$, donde $E\in\Ob D(V)$, y que se tiene un isomorfismo
\[
\alpha:f_V^*E\xrightarrow{\sim}\Lambda_X\lten_{\Lambda}E_0
\]
de $D(G,X)$, con $E_0\in\Ob D(\Lambda[G])$, perfecto sobre $\Lambda$. Sea
\[
v\in\Hom(d_1^*M,d_2^*M)
\]
como en 6.21, y denotemos por
\[
v_0\in\Hom_{\Lambda[G]}(E_0,E_0)
\]
el homomorfismo definido, gracias a $\alpha$, por $f^*v|c_1^{-1}(U)$. Entonces, si $\underline c_{U,X}$ designa la correspondencia definida en 6.14.3, se tiene, con la notación 6.20.2,
\begin{equation}
\Tr_{\Lambda}(v^!)
=
\sum_{g\in G_{\natural}}
\Tr_{\Lambda[G]}(f_*\underline c_{U,X})(g^{-1})\Tr_{\Lambda}(gv_0).
\tag{6.23.1}
\end{equation}
\end{statement}

Obsérvese que, gracias al carácter central de $\Tr_{\Lambda}$, el segundo miembro no depende de la elección de $\alpha$.

Por la fórmula de proyección
\[
f_*f^*\Lambda_Y\lten M\xrightarrow{\sim}f_*f^*M
\]
y 6.15, $M$ satisface las hipótesis de 6.22. Por otra parte, $\alpha$ define un isomorfismo
\[
f^*M\xrightarrow{\sim}\Lambda_{U,X}\otimes E_0
\]
mediante el cual $(f^*v)^!$ se identifica con el producto tensorial externo $\underline c_{U,X}\otimes v_0$. Por la fórmula de proyección, $f_*f^*M$ se identifica, pues, con $f_*(\Lambda_{U,X})\otimes E_0$, y $f_*((f^*v)^!)$ con el producto tensorial externo $f_*\underline c_{U,X}\otimes v_0$. Para $g\in G$, la correspondencia $Z_g$-equivariante $gf_*((f^*v)^!)$ se identifica, por tanto, con
\[
gf_*\underline c_{U,X}\otimes gv_0.
\]
Como, según 6.10.4, se tiene
\[
\Tr_{\Lambda[G]}(f_*((f^*v)^!))(g^{-1})
=
\Tr_{\Lambda[Z_g]}(gf_*((f^*v)^!))(e),
\]
basta, teniendo en cuenta 6.22.1, demostrar la fórmula
\[
\Tr_{\Lambda[Z_g]}(gf_*\underline c_{U,X}\otimes gv_0)(e)
=
\Tr_{\Lambda[G]}(f_*\underline c_{U,X})(g^{-1})\Tr_{\Lambda}(gv_0),
\]
o, lo que viene a ser lo mismo según 6.10.4,
\begin{equation}
\Tr_{\Lambda[Z_g]}(gf_*\underline c_{U,X}\otimes gv_0)(e)
=
\Tr_{\Lambda[Z_g]}(gf_*\underline c_{U,X})(e)\Tr_{\Lambda}(gv_0).
\tag{6.23.2}
\end{equation}
Sustituyendo, si es necesario, $G$ por $Z_g$, se puede suponer que $g=e$ y, volviendo a la definición de las trazas, se ve que basta demostrar el lema siguiente.
